\section{Conclusion}
The methodology detailed in this chapter, from the physics-aware data pre-processing to the containerized, modular deep learning implementation, establish a robust foundation for the experiments.
By automating data ingestion, ensuring memory safety through chunking, and rigorously isolating the execution environment via Docker, we have minimized computational overhead and procedural errors and maximized the scientific validity of the results.
Furthermore, the integration of the Optuna optimization framework enables a systematic exploration of the model's parameters space, ensuring that the configurations used in the subsequent analysis are data-driven rather than arbitrary.
The following chapter will present the quantitative and qualitative outcomes of our experiments, demonstrating the effectiveness of the proposed CycleGAN approach in translating synthetic images into realistic observations for the ultimate goal of unsupervised subsurface anomaly detection.

